\documentclass[11pt,a4paper]{article}
\usepackage{amsmath,amssymb,amsthm,mathtools,booktabs,microtype}
\usepackage[margin=25mm]{geometry}
\usepackage[colorlinks=true,linkcolor=blue,urlcolor=blue]{hyperref}
\setlength{\emergencystretch}{2em}
\newtheorem{theorem}{Theorem}[section]
\newtheorem{lemma}[theorem]{Lemma}
\newtheorem{proposition}[theorem]{Proposition}
\newcommand{\Q}{\mathbb Q}\newcommand{\Z}{\mathbb Z}
\DeclareMathOperator{\Tr}{Tr}\DeclareMathOperator{\Nm}{Nm}
\title{One-helper pole normalization:\\an ordinary-abc quartic in 58 coordinates}
\author{\(\Z\Q\) Diophantine Project}
\date{19 November 2026}
\begin{document}\maketitle
\begin{abstract}
We reduce the audited 59-coordinate ordinary-abc quartic to 58 coordinates.
The one-helper normalization $\alpha=u+u^2D^2$ and
$b=7u^{-1}(\alpha/(1+\alpha))^2$ simultaneously preserves extraction at
positive odd valuations and controls every negative odd valuation,
without assuming that auxiliary-point denominators were fixed before prime
selection. Its only division is guarded by its own quadratic equation.
The factor seven puts every prepared pole input in the square class required
by Burungale's published simultaneous-twist theorem, without a further
square-class extension. The result is 50 integral quadrics in 58 rational
coordinates, whose sum of squares has 867 monomials, degree four, and
integer content one. Every sufficiently large fixed odd shift defines the
integers under ordinary rational abc and the explicitly retained elliptic,
comparison, trace-index, and prime-pattern premises. The sufficient shift
remains ineffective.
\end{abstract}

\section{Statement, attribution, and retained inputs}

The shared detector, translated-secant chart, reference rechart, and
Laurent packing used below originate in the notes of \textbf{Gabe Conant
with Astra 6.2}, supplied as \texttt{H10-quartic-constructions.zip} and
\texttt{H10-eleven-variable-constructions.zip}. Their 62-coordinate
construction was audited and extended in v78~\cite{Conant,v78}.
The single-pole derivative transfer and compressed derivative circuit were
proved in v79~\cite{v79}. The one-helper normalization here is subsequent
project work. Its pole-completeness input is derived below directly from
the stated square classes in Burungale~\cite[Theorem B]{Burungale}.

We retain the parent assertions concerning the elliptic group over
$K=\Q(\sqrt2)$, the inverse-parameter extraction theorem and its formal
comparison estimates, the trace-index bound for all allowed points, the
simultaneous prime-pattern theorem with finite avoidance and local and
real preparation. Ordinary rational abc supplies both the derivative
contact estimate and the signed mixed-power estimate. These retained
premises are not claimed to follow from abc. For the native pole we use
Burungale's Theorem B with its published congruences~\cite{Burungale};
no assertion for the square class $1$ replaces its required class $7$.

\begin{theorem}\label{thm:main}
Under these retained premises and ordinary rational abc, there is an
ineffective constant $B_0$ such that, for every odd integer $B\ge B_0$,
the polynomial specified below has rational input projection $\Z$.
It has 58 total variables, including the input $X$, and is the sum of
50 squares of integral quadrics. Its degree is exactly four, its content
is one, and it has 867 coordinate monomials. For every positive integer
$B$, its exact coefficient height is
\[
 H(F_B)=\max\{22838305425411,\ B^2+4091\}.
\]
This coefficient formula is algebraic and is not a sufficient defining
threshold for $B$.
\end{theorem}

The complete expanded coefficient list and quadratic specification are
\nolinkurl{computation/variable_reductions_v80/abc/quartic58/polynomial.json}
and \nolinkurl{quadrics.json} in that directory. A finite reconstruction
does not certify the retained arithmetic premises.

\section{The unchanged elliptic chart}

Put $\rho^2=2$, $\Delta=11-6\rho$, and
\[
 E:y^2=x^3-\Delta x,\qquad R_0=(1+\rho,2\rho).
\]
We use exactly the complete paired chart of v65 as recharted in
v78. For $X\ne0$ it gives finite points $U,V,H\in E(K)$ with the
specified nonzero abscissas and an allowed reference, with $U\ne V$.
It includes tangencies, zero ordinates, and every finite two-torsion
difference. Its physical outputs are
\[
\begin{gathered}
 s=X\,y(U)/x(U),\qquad t=X\,y(V)/x(V),\\
 S=y(H)/x(H),\qquad T=y(H-R_0)/x(H-R_0),\\
 J=X^2f(H)f(U),\qquad K'=X^2f(H-R_0)f(V),\qquad f=1/x.
\end{gathered}
\]
Let $C=s-S$ and $I=t-T$. The stored rational half-trace $h_0$ satisfies
\[
 T_{\mathrm{sc}}=2h_0=X^2\Tr x(U-V).
\]
The retained index $m\in\Z$ of $U-V$ has the parent trace bound
\begin{equation}\label{eq:index}
 m^2\le C_0\bigl(h(T_{\mathrm{sc}})+2h(X)\bigr)+J_0,
\end{equation}
with fixed constants. Finite torsion differences are included with
$m=0$. The parent inverse-parameter extraction theorem says that
positivity of the four valuations of $C,J,\overline I,\overline{K'}$
at one good slot implies that $X$ is locally integral and
\begin{equation}\label{eq:local-label}
                         X(m-X)\equiv0\pmod p.
\end{equation}
All fixed exceptions in these statements remain excluded.

For a nonzero rational $c$ define
\begin{equation}\label{eq:detector}
 W=1+cC^4+c^2J^2-c^3\overline I^{\,4}-c^4\overline{K'}^{\,2},
 \qquad D\Tr W=\Nm(W-1).
\end{equation}
Every rational solution has $\Tr W\ne0$: otherwise the norm vanishes,
so $W=1$ in the field $K$, contrary to $\Tr 1=2$.
The even valuations of reciprocal abscissas give even valuations of
$J,K'$ at every good unramified slot, with zero values interpreted as
having infinite valuation.

\section{The contact quotient and its automatic guards}

Set
\[
 M=B+X^8+T_{\mathrm{sc}}^2,\qquad a=M^4,\qquad
 x=a-z,\quad y=a+z,
\]
where $z$ is the one free contact parameter. Let
\[
 \mathcal C=\{-8,-7,-2,-1,3,4,5,6\},\qquad
 G(t)=\prod_{r\in\mathcal C}(t+r),
\]
\begin{equation}\label{eq:derivative}
 F=G(x)G(y),\qquad F_a=G'(x)G(y)+G(x)G'(y),\qquad
 u=-16F/F_a,\qquad c=1/u.
\end{equation}
The shifts are distinct, nonzero, have sum zero, and no two are
opposites. In particular $2a+r+s>0$ for all shifts when $B\ge3$.

The equations used in the quadratic chart are $c c^{-1}=1$ and
$F_a/c+16F=0$, together with the exact value and derivative evaluations.
They force $F F_a\ne0$ for $B\ge3$. Indeed $G(x)=G(y)=0$ would give
$x+y=-(r+s)$ of absolute value at most 16, whereas $x+y=2M^4>16$.
If just one factor is zero, its derivative is nonzero because $G$ is
squarefree, and the other factor is nonzero. Thus $F_a\ne0$, contradicting
the displayed relation. If $F_a=0$, that relation first gives $F=0$,
which was just excluded. No inverse of $F$ or $F_a$ needs another
coordinate.

The intermediate derivative-height theorem of v63~\cite{v63}, before
replacing the larger slot height by half the product height, gives for
$\epsilon=1/100$
\[
 \log r^+(u)\ge\frac{197}{100}\max\{h(x),h(y)\}
                         -\frac{201}{100}h(M)-O(1),
\]
where $r^+(u)$ is the product of primes of positive odd valuation
outside a fixed set. The fixed set includes every bad chart prime,
shift difference, nonzero pair sum, and the primes of $16\cdot7$.
The constants are independent of $B$. Since $x+y=2M^4$,
\begin{equation}\label{eq:radical}
             \log r^+(u)\ge\frac{193}{100}h(M)-O(1).
\end{equation}
The proof of the intermediate estimate remains valid if $x$ or $y$
vanishes: its eight shifted contact factors are then nonzero constants,
and the larger-height slot supplies the simple-prime argument.
There is no division by $xy$ in this construction.

Because $M\ge B$, estimate~\eqref{eq:radical} implies that $u$ has
at least one positive odd valuation for every allowed rational tuple once
$B$ is sufficiently large. In particular $u$ is not a negative rational
square. This consequence is independent of the pole equations and will
exclude a possible zero of the new normalizer.

\section{One helper, both signs of odd valuation}

Define
\begin{equation}\label{eq:normalizer}
 \alpha=u+u^2D^2,\qquad
 \tau=\frac{\alpha}{1+\alpha},\qquad
 b=7c\tau^2=\frac7u\left(\frac{\alpha}{1+\alpha}\right)^2.
\end{equation}
The gate is the polynomial equation $\tau(1+\alpha)=\alpha$.
It excludes $\alpha=-1$ identically. If $\alpha=0$, then
$u=-(uD)^2$, contrary to~\eqref{eq:radical} for sufficiently large $B$.
Thus $\tau$ and $b$ are nonzero on every sufficiently large defining fiber.
No additional nonzero coordinate is needed.

The native pole equations are
\begin{equation}\label{eq:pole}
 e^2+bp^2=1,\qquad u_0v_0=e(1-b),\qquad u_0^2-v_0^2=bq^2.
\end{equation}
Here $u$ is the contact quotient; $u_0,v_0$ are pole coordinates.

\begin{lemma}\label{lem:pole-obstruction}
If a nonzero rational $b$ admits a rational solution of~\eqref{eq:pole},
then $v_p(b)$ is never a negative odd integer, for any prime $p$.
No nonzero hypothesis on $e,p,u_0,v_0,q$ is needed.
\end{lemma}
\begin{proof}
Suppose $v_p(b)<0$ is odd. The two terms on the left of the first
equation have distinct valuation parity. They cannot cancel, so both
are integral. The term $bp^2$ cannot have valuation zero; hence $e$
is a unit. The middle equation then gives
$v_p(u_0)+v_p(v_0)=v_p(b)$, so both factors are nonzero and their
valuations are unequal. Thus $u_0^2-v_0^2$ has even valuation, whereas
$bq^2$ has odd valuation if $q\ne0$ and is zero if $q=0$.
Both alternatives are impossible.
\end{proof}

\begin{lemma}\label{lem:normalizer}
Let $p\ne7$, let $\mu=v_p(u)$ be odd, and put $d=v_p(D)$, allowing
$d=+\infty$. Then
\[
 A:=v_p(\alpha)=\min\{\mu,2\mu+2d\}.
\]
If $\mu<0$, then $v_p(\tau)=0$ and $v_p(b)=-\mu>0$.
If $\mu>0$ and the native pole equations hold, then
$A>\mu/2$ and $d>-3\mu/4$.
\end{lemma}
\begin{proof}
The two summands defining $\alpha$ have different valuation parity;
hence their minimum is attained once. In particular $\alpha\ne0$.
For $\mu<0$ we have $A\le\mu<0$, so $1+\alpha$ has valuation $A$
and $\tau$ is a unit. This proves the first conclusion for arbitrary $D$.

Suppose $\mu>0$. If $A<0$, then $v_p(\tau)=0$ and $v_p(b)=-\mu$,
contradicting Lemma~\ref{lem:pole-obstruction}. If $A=0$, the gate
ensures its denominator is nonzero, with $v_p(1+\alpha)\ge0$.
Thus $v_p(\tau)\le0$, again giving a negative odd valuation of $b$.
Consequently $A>0$, so $v_p(\tau)=A$ and
$v_p(b)=-\mu+2A$ is odd. It must be positive. Therefore
$A>\mu/2$, which forces $2\mu+2d>\mu/2$, as claimed.
\end{proof}

The negative-odd conclusion is uniform in the actual point denominators.
No second round of local preparation is needed after the auxiliary points
are chosen. The prime seven is included in the fixed exceptional set and
is treated explicitly in the completeness preparation.

\section{Soundness}

Take a consumed prime with $\mu=v_p(u)>0$ odd. Then $v_p(c)=-\mu$.
If both slot valuations of $W$ are negative, say $w_1,w_2$, then
\[
 v_p\Nm(W-1)=w_1+w_2,\qquad
 v_p\Tr W\ge\min(w_1,w_2),
\]
so $v_p(D)\le\max(w_1,w_2)$. Otherwise a slot is already integral.
Lemma~\ref{lem:normalizer} therefore gives a slot at which
$v(W)>-3\mu/4$. This includes an inert prime with repeated slot
valuations.

At that slot the four nonconstant terms of~\eqref{eq:detector} have
pairwise distinct classes modulo four, namely $-\mu,-2\mu,-3\mu,0$.
Here the second and fourth powers have the required classes because
$v(J),v(K')$ are even. Every negative minimum is unique. It follows
that each nonconstant term has valuation greater than $-3\mu/4$.
In particular
\[
 v(C)>\mu/16,\quad v(J)>5\mu/8,\quad
 v(\overline I)>9\mu/16,\quad v(\overline{K'})>13\mu/8.
\]
All four tests are positive. Zero tests satisfy the inequalities with
infinite valuation. The retained extraction theorem now gives
\eqref{eq:local-label}.

If $X\ne0$ and $X\ne m$, writing $X=A/Q$ in lowest terms gives
\[
 \log r^+(u)\le\log|A(mQ-A)|
             \le2h(X)+\log(1+|m|).
\]
Apply the signed mixed-power estimate~\cite{v58} with exponent eight
and $\kappa=1001/1000$ to $s=X^8+T_{\mathrm{sc}}^2$. Since
$h(s)\le h(M)$ for $B>0$,
\[
 h(X)\le\frac{501}{997}h(M)+O(1).
\]
Combining with~\eqref{eq:radical} leaves the strict positive margin
\[
 \frac{193}{100}-\frac{1002}{997}
                       =\frac{92221}{99700}>0.
\]
It follows that $h(M)$ is bounded by a fixed multiple of
$\log(1+|m|)+1$. Also
\[
 2h(T_{\mathrm{sc}})\le h(s)+8h(X)+\log2.
\]
The trace estimate~\eqref{eq:index} then bounds $m$, and therefore
$h(M)$, independently of $B$. This contradicts $h(M)\ge\log B$
for sufficiently large $B$. Thus every rational input solution has
$X=0$ or $X=m\in\Z$.

\section{Complete integer witnesses}

Fix $X=n\ne0$. Use the parent multiplication family with
$U-V=[n]R_0$, so $T_{\mathrm{sc}}$, $M$ and $a=M^4$ are fixed before
contact prime selection. The sixteen rational positions $\pm(a+r)$,
$r\in\mathcal C$, are distinct. Apply the retained simultaneous
prime-pattern theorem with all fixed exceptions, bounded real and local
point neighborhoods, and the same compatible comparison-index refinements
as in the parent construction. These are sixteen simple contacts instead
of the two previous eight-contact lists; the retained theorem allows any
fixed finite list. At each selected contact prime exactly one of the
sixteen linear factors vanishes simply, its opposite factor is a unit,
and therefore $v_p(u)=1$ by the derivative identity.

The parent successful-slot preparation gives $v(C),v(\overline I)\ge3$
and $v(J),v(\overline{K'})\ge4$ there. Thus $W\equiv1$. At the other
slot the prepared reciprocal products are units and the fourth term in
\eqref{eq:detector} has the unique minimum $-4$. The shifted norm and
trace give $v_p(D)>0$, including $D=0$. Thus $v_p(\alpha)=1$,
$v_p(\tau)=1$, and $v_p(b)=1$ at every selected prime. The selected
primes avoid seven.

For this fixed integer witness, enlarge the finite preparation set to
$S_n$ by the coefficient and denominator primes of the now-fixed contact
positions, and by the primes of the nonzero rational $a$. This does not
enlarge the fixed exceptional set in the height theorem. Choose a fixed
positive integer $D_0$ with $D_0^2(a+r)\in\Z$ for every shift and take
the prepared slope in the form
\[
 z=\frac{N}{D_0^2L^2},\qquad
 p_{r,\pm}=N\pm D_0^2(a+r)L^2.
\]
The retained prime supplier gives the sixteen distinct positive primes
$p_{r,\pm}$ with the required finite avoidance. Exactly,
\begin{equation}\label{eq:prime-factorization}
 F=\frac{\prod_{r\in\mathcal C}p_{r,-}p_{r,+}}
                {D_0^{32}L^{32}}.
\end{equation}
The sign is positive because there are eight quadratic factors.

Every other prime is covered by the following three cases. Outside
$S_n$, if $z$ is integral and the prime is not one of the sixteen selected
primes, its denominator is a unit: otherwise that prime divides $N,L$
and hence all sixteen distinct selected primes, a contradiction.
Formula~\eqref{eq:prime-factorization} then gives $v_p(F)=0$.
The derivative is integral, so $v_p(u)\le0$. If $v_p(z)<0$ outside
$S_n$, the leading terms are $F=z^{16}+O(z^{14})$ and
$F_a=-16az^{14}+O(z^{12})$ with $a$ a unit. Their strict dominant
valuations give $v_p(u)=2v_p(z)$, an even integer. At the finitely many
places in $S_n$, including any numerator or denominator prime of $a$,
the prescribed sufficiently deep slope denominator gives instead
\[
                      v_p(u)=2v_p(z)-v_p(a).
\]
This is again even because $a=M^4$. The coefficient 16 cancels between
the quotient numerator and the derivative leading term, including at
two. Thus all positive odd valuations of $u$ occur among the selected
primes already checked.

At any negative odd valuation of $u$ outside seven, including an
uncontrolled odd factor of the derivative denominator,
Lemma~\ref{lem:normalizer} applies without any restriction on the actual
point coordinates and gives positive odd valuation of $b$.
If $p\ne7$ and $v_p(u)$ is even, then $v_p(b)=-v_p(u)+2v_p(\tau)$
is even. The exceptional prime seven is prepared below with positive
valuation of $b$. Consequently the actual $b$ has square reduced
denominator once that fixed preparation is imposed.

At fixed places and real infinity in the slope, the zero shift sum gives
\[
 u\sim z^2/a,\qquad c\sim a/z^2.
\]
The already fixed point neighborhoods keep the raw tests bounded.
Consequently $W-1=O(c)$, $D=O(c^2)$, and
\[
 \alpha=c^{-1}+O(c^2),\qquad \tau=1+O(c),\qquad b\sim7c.
\]
At the real place take $z$ sufficiently large; then $c>0$, $\alpha>0$,
$0<\tau<1$, and $0<b<1$. At two, $a=M^4$ is a rational square.
Prescribing an arbitrarily deep slope denominator makes $c$ a two-adic
square of arbitrarily large positive even valuation, while $\tau$ is
a nonzero unit. Therefore
\[
 b\in7\Q_2^{\times2},\qquad v_2(b)\ge4.
\]
At every other fixed place, the same preparation makes $v_p(c)$ positive
and even, and $\tau$ a unit. In particular $v_7(b)=1+v_7(c)>0$.
These preparations precede selected prime choices and impose only the
permitted finitely many local conditions.

The next lemma therefore applies to the actual detector value $b$.

\begin{lemma}\label{lem:pole-completeness}
Suppose that $0<b<1$, that $b$ has square reduced denominator, and that
$b\in7\Q_2^{\times2}$ with $v_2(b)\ge4$. Then~\eqref{eq:pole} has a
rational solution.
\end{lemma}
\begin{proof}
Write $b=N/Q^2$ in lowest terms with $0<N<Q^2$ and $Q>0$.
The two-adic assumptions make $Q$ odd and $16\mid N$.
Let $\beta,\eta$ be the positive squarefree parts of $N,Q^2-N$,
respectively. Thus
\[
 b=\beta z_0^2,\qquad 1-b=\eta d_0^2
\]
for nonzero rational $z_0,d_0$. We have $\gcd(\beta,\eta)=1$,
$\beta\equiv7\pmod8$, and $\eta\equiv1\pmod8$.
If $p\mid\beta$, then $1-b\equiv1\pmod p$, so $\eta$ is a square
in $\Q_p$. If $p\mid\eta$, then $b\equiv1\pmod p$, so $\beta$
is a square in $\Q_p$. These are odd primes, so these residue statements
lift to the required local squares. The value $\eta=1$ is allowed.

Apply precisely Theorem B of Burungale~\cite{Burungale}. It supplies
an odd positive squarefree $n\equiv5\pmod8$ and points on
$E_n$ and $E_{n\eta}$ whose nonzero abscissas lie in
$\beta\Q^{\times2}$, where $E_l:y^2=x^3-l^2x$.
Their ordinates are nonzero: the positive nonzero two-torsion abscissas
are $n$ and $n\eta$, with two-adic square class $5$, whereas the
required class is $\beta\equiv7\pmod8$. Negative and zero
abscissas are excluded by $\beta\Q^{\times2}$.
Scale the second point by $(x,y)\mapsto(d_0^2x,d_0^3y)$.
We obtain points
\[
 (ba_1^2,y_1)\in E_n(\Q),\qquad
 (ba_2^2,y_2)\in E_{n(1-b)}(\Q),\qquad a_1a_2\ne0.
\]
Define
\[
 e=\frac{n}{ba_1^2},\qquad p=\frac{y_1}{b^2a_1^3},\qquad
 u_0=\frac{a_2}{a_1},\qquad
 v_0=\frac{e(1-b)}{u_0},\qquad
 q=\frac{y_2}{b^2a_1^3u_0^2}.
\]
All denominators are nonzero. The first curve equation gives
$e^2+bp^2=1$; the definition of $v_0$ gives the middle pole equation;
and the second curve equation gives $u_0^2-v_0^2=bq^2$.
This proves all three equations for this exact rational input.
\end{proof}

For $X=0$, retain the checked parent zero fiber: $H=2R_0$,
$s=S$, $t=T$, $F_{\mathrm{pair}}=G_{\mathrm{pair}}=d_{
\mathrm{pair}}=J=K'=0$, and
$h=t^2/(s-t)$, whose denominator is nonzero because
$\Nm(S-T)=7199/328$. Its scaled trace fixes $M$ before the contact
choice. Every raw test vanishes, so $W=1$, $D=0$, $\alpha=u$,
$\tau=u/(1+u)$, and $b=7c/(1+c)^2$. The same contact, real and
fixed-place preparations and Lemma~\ref{lem:pole-completeness} supply
a rational zero. No input factor is
multiplied into the quartic.

\section{A quadratic contact circuit without cubic-power registers}

Write $G(t)=\sum_{i=0}^8a_it^i$. Our shifts give $a_8=1$, $a_7=0$,
and $a_6=-102$. Define
\[
 P(t)=tG'(t)-8G(t),\qquad
 Q(t)=(4t^2+a_6)G'(t)-32tG(t).
\]
Both have degree at most six. The degree-nine terms in the second
identity cancel; its degree-eight term vanishes because $a_7=0$, and
its degree-seven coefficient is $-8a_6+8a_6=0$.

\begin{lemma}\label{lem:derivative-circuit}
Store $t,t_2,t_4,G_t,H_t$ with $t_2=t^2$, $t_4=t_2^2$, the exact
quadratic evaluation $G_t=G(t)$, and
\[
 tH_t-8G_t=P(t),\qquad
 (4t_2+a_6)H_t-32tG_t=Q(t).
\]
Every displayed residual is quadratic in these five coordinates, and
the equations imply $H_t=G'(t)$ at every rational point, including $t=0$.
\end{lemma}
\begin{proof}
Every power through degree six is a product of at most two among
$1,t,t_2,t_4$, and $t^8=t_4^2$; $G$ has no degree-seven term.
Subtract the exact derivative identities from the stored equations.
They give $t(H_t-G')=0$ and $(4t^2+a_6)(H_t-G')=0$.
The identity $(4t^2+a_6)-4t\,t=a_6=-102\ne0$ proves the claim.
The second equation supplies exactly the missing zero-coordinate fiber.
\end{proof}

There are two such evaluations, one for each slot. The sum
$F_a=G'(x)G(y)+G(x)G'(y)$ needs one further stored coordinate.
Together they use two fewer coordinates than the direct circuit with
separate $x^3,y^3$ registers. The extra derivative rows are algebraically
dependent on the nonzero slots; row count is not being used as a
dimension or independence assertion.

\section{Exact compilation and coefficient verification}

Use the original v78 Laurent-scaled coordinates for the reference,
paired points, hyperbolic test pairs, input powers, $W,D$, and the first
native pole. Store $c,c^{-1}$ with their literal unit equation.
In particular the physical-to-stored conventions are
\[
 D=cD_s,\quad b=cb_s,\quad e=ce_s,\quad
 u_0=cu_s,\quad v_0=cv_s,\quad p=p_s,\quad q=q_s,
\]
\[
 p^2=cP_s,\qquad q^2=cQ_s.
\]
The five unchanged pole rows therefore read
\[
 p_s^2=cP_s,\quad e_s^2+b_sP_s=c^{-2},\quad
 u_sv_s=e_sc^{-1}-b_se_s,
\]
\[
 q_s^2=cQ_s,\qquad u_s^2-v_s^2=b_sQ_s.
\]
These equations decode exactly to~\eqref{eq:pole}; in particular
there is no hidden additional square in the physical input $b$.

The only new normalizer coordinate is $\alpha$; $\tau$ is retained
from the parent pole input. The three new rows are
\begin{equation}\label{eq:quadratic-gate}
 \alpha=c^{-1}+D_s^2,\qquad
 \tau(1+\alpha)=\alpha,\qquad b_s=7\tau^2.
\end{equation}
Since $D=cD_s$ and $u=c^{-1}$, they recover
$\alpha=u+u^2D^2$ and $b=7c\tau^2$ exactly.
The middle row forces $1+\alpha\ne0$ without any additional guard.
At large defining shifts, $\alpha\ne0$ follows separately from the
derivative-height estimate, as proved above.

Store $M,M_2$ unscaled with
\[
 M=B+X_4^2+4h_0^2,\qquad M_2=M^2.
\]
Here $X_4=X^4$ is the retained input-power coordinate. For both contact
slots, physical $t,t_2,t_4,G_t,H_t$ are $c$ times their stored
coordinates, while physical $F_a$ is $c^2$ times its stored coordinate.
The midpoint equation is $c(x_s+y_s)=2M_2^2$.
Divide each contact evaluation and derivative row by $c^2$;
Lemma~\ref{lem:derivative-circuit} shows that the resulting residuals
remain quadratic, with $c^{-1}$ and $c^{-2}=(c^{-1})^2$ accounting for
linear and constant terms. The last two rows are
\[
 (F_a)_s=(G'_x)_s(G_y)_s+(G_x)_s(G'_y)_s,
 \qquad (F_a)_sc^{-1}+16(G_x)_s(G_y)_s=0.
\]

\begin{samepage}
The exact census is
\begin{center}
\begin{tabular}{lr}
\toprule
Coordinates&Count\\\midrule
Input, reference and paired chart, four tests, input powers&30\\
Shared $W$ and $D$&3\\
One native pole, its input, and $\tau$&9\\
$c,c^{-1}$&2\\
$M,M_2$&2\\
Two contact slots, each with two powers, value and derivative&10\\
Derivative sum&1\\
One-helper normalizer $\alpha$&1\\\midrule
Total&58\\\bottomrule
\end{tabular}
\end{center}
\end{samepage}
The complete list has 50 quadrics. Their sum of squares vanishes over
$\Q$ precisely when every residual vanishes. Every change of coordinates
used above has a recovered nonzero denominator, including the native
zero-coordinate cases.

Direct exact expansion gives 867 coordinate monomials and 870
parameter-separated terms. The maximum coefficient independent of $B$
is $22838305425411$. The only mixed profiles are
$B^2+4091$, $8B+1$, and $2B+1$; the remaining parameter-dependent
coefficients are bounded by these or by the fixed maximum. For $B>0$,
$B^2+4091$ dominates the two linear mixed profiles and all pure linear
profiles. This proves the height formula in Theorem~\ref{thm:main}.
The coefficient of $X^4$ is one, proving both exact degree four and
content one. All 58 coordinates occur.

The independent checker reconstructs the retained chart from its pinned
parent, decodes every new physical scaling, derives the two derivative
identities independently, expands every squared row with independent
SymPy polynomial arithmetic, and checks the census, support, content,
height profiles and exact rational height margin. The universal
valuation and completeness arguments are the proofs above, not an
inference from a finite residue sample.

\section{Scope of the improvement}

This is a reduction from 59 to 58 coordinates for the ordinary-abc
quartic route. The prime-seven pole input also restricts the twist input to the exact
square classes stated in the published theorem. The normalizer is reusable wherever a rational $u$ is
required to detect positive odd valuations while its negative odd
valuations must remain admissible for a native pole. It does not supply
a rational one-parameter specialization of the intrinsic pole surface.
Accordingly, this paper does not reduce the separate eleven-coordinate
unbounded-degree construction or resolve its proposed nine-coordinate
actual-realization condition. No global minimum is asserted.

\begin{thebibliography}{9}\footnotesize\raggedright
\setlength{\itemsep}{1pt}
\bibitem{Conant} G.~Conant, with Astra 6.2,
supplied notes \texttt{H10-quartic-constructions.zip} and
\texttt{H10-eleven-variable-constructions.zip}, November 2026;
particularly \texttt{abc62/PROOF.md} and \texttt{fewvars/BALANCED.md}.
The original provenance is retained in the repository's v78 package.
\bibitem{v79} \(\Z\Q\) Diophantine Project,
\emph{One derivative contact and one pole: an ordinary-abc quartic in 59 coordinates},
\nolinkurl{papers/quartic_existence/single_derivative59_20261119_v79.tex}.
\bibitem{Burungale} A.~Burungale,
\emph{A Diophantine definition of $\Z$ in $\Q$},
arXiv:2611.16130v1, November 2026, Theorem B,
\url{https://arxiv.org/html/2611.16130v1}.
\bibitem{v78} \(\Z\Q\) Diophantine Project,
\emph{Audit and extensions of the 62-coordinate quartic},
\nolinkurl{papers/quartic_existence/abc_quartic_audit_20261119_v78.tex}.
\bibitem{v65} J.~Brox and contributors,
\emph{A quartic in 68 variables from translated secants},
\nolinkurl{papers/quartic_existence/quartic_abc68_20261117_v65.tex}.
\bibitem{v63} \(\Z\Q\) Diophantine Project,
\emph{An eleven-variable derivative detector},
\nolinkurl{papers/conditional_few_variables/derivative_detector11_20261117_v63.tex}.
\bibitem{v58} \(\Z\Q\) Diophantine Project, signed mixed-power height theorem,
\nolinkurl{papers/elliptic_heights/signed_mixed_power_height_20261117_v58.tex}.
\end{thebibliography}
\end{document}
